% ======================================================================
% ABSTRACT (title page block, single paragraph approximately 1000 chars)
% Location: 2030-pdac-1min/paper/full-paper/sections/abstract.tex
% ======================================================================

On premises repository based large language models (LLMs) provide commands to standard oncology surgical robots based on real time sensor data and are controlled via x, y, z coordinates to administer patient treatment, which minimizes single robot error potential. The thesis is operationalized through a four phase Claude Code workflow (instructions plus codegen plus execution plus draft paper) producing a 60 second 8 arm robotic pancreaticoduodenectomy on a hypothetical 2030 Medtronic PancreSpeed 1.0 platform paired with the Daraxonrasib pan KRAS inhibitor (FDA Breakthrough Therapy 06/25, RASolute 302 + RASolve 301). PancreSpeed 1.0 closes a 5x to 500x SOTA gap (1,200 mm/s tip, 100 kHz force, 0.05 mm RMS, 3 ms E stop, 18 N cumulative cap, 640 sensor channels). A 32 iteration deterministic Latin hypercube sweep at seed 20260513 yields mean composite 93.298 (std 1.225, 95\% CI half width 0.462), and the 4 entrant LLM tournament places PancreSpeed 1.0 first at 93.735 versus 2030 da Vinci SP and Hugo RAS successors and the 2025 Dutch human baseline (47.0\% ideal). The 1001 record Phase 5 first 100 ms \texttt{sensor\_sample\_8arm.jsonl} is the exceptional processing feat; Daraxonrasib postop restart distributes as T+7d in 29 of 32 iterations. Practical adoption still requires bridging the synthetic patient gap, hypothetical hardware gap, and TRIPOD+AI plus CREMLS reporting floors (\S 1.4, \S 4.5).
